\subsection{The politico-economic equilibrium with CRRA}\label{\refsec:PEE}
With CRRA preferences the political problem carries \emph{two} endogenous states. The informal savings ratio $\iota_{t-1}$ is a state for the reason given in \cref{\refsec:PEELOG}, and it survives unchanged. The savings level $s_{t-1}$ returns as a second state because the argument of \eqref{\refeq:logsep} no longer applies: the marginal indirect utilities are now (level)$^{1-1/\rho}\times$(log-derivative) rather than pure log-derivatives, and while the log-derivatives remain independent of $s_{t-1}$, the levels do not. This subsection sets out how to search over the pair without the cost, or the silent failure modes, that a naive product grid invites.

\smalltitle{Grids.}
Retain $\mathcal{T}$, $\mathcal{S}_0$ and $\mathcal{S}_0'$ from \cref{\refsec:PEELOG}, including the restriction $l_{\iota}>0$ and its justification. Add $\mathcal{S} = \left\lbrace s^{(1)},...,s^{(M_s)}\right\rbrace$ with $s^{(1)} = l_s>0$ and $s^{(M_s)} = u_s>l_s$, a grid of predetermined savings levels, and $\mathcal{S}'$, a separate grid of \emph{candidate} values of $s_t$ -- the same distinction as between $\mathcal{S}_0$ and $\mathcal{S}_0'$. Realistic bounds for $\mathcal{S}$ come from the steady state, but the choice of \emph{where} to evaluate it matters more than it appears. Since $s^*(\tau)$ is monotonically decreasing in $\tau$ with $s^*\rightarrow 0$ as $\tau\rightarrow 1$, the obvious anchor is $s^*(0)$, and $u_s = 1.25\times s^*(0)$ with $l_s$ a small positive value is serviceable at any single calibration. It does not survive a sweep. Across $\rho\in[0.5,2]$, $s^*(0)$ \emph{falls} by $83\%$ while the upper edge of the reachable set \eqref{\refeq:reachable} \emph{rises} by $20\%$: the two are almost perfectly negatively correlated, so no fixed multiple of $s^*(0)$ tracks the region the dynamics occupy, and the fraction of $\mathcal{S}$'s nodes that lie inside it swings from $0.4$ to $0.8$ across the sweep. Evaluating at the calibration's own target $\tau_0$ is better but still drifts by $77\%$. The difficulty is that $s^*(\tau)$ at low $\tau$ is dominated by the same crowding-out channel that makes $\iota^*$ diverge at the other end, and so moves with $\rho$ for reasons unrelated to where savings actually settle. At an interior $\tau$ -- we use $\tau=0.3$ -- the steady state is stable to $1.5\%$ across the same range while remaining a solved function of the calibrated parameters, which is what an anchor has to be: responsive to the data, not to the parameter being swept. We therefore take $l_s = \max\lbrace\delta_s,\,\underline{\lambda}_s s^*(\bar\tau)\rbrace$ and $u_s = \overline{\lambda}_s s^*(\bar\tau)$ with $\bar\tau = 0.3$, $\underline{\lambda}_s = 0.45$ and $\overline{\lambda}_s = 3.65$, again carrying about $20\%$ margin on the measured edges. The state grid is $\mathcal{S}\times\mathcal{S}_0$ and the search grid is $\mathcal{T}\times\mathcal{S}'\times\mathcal{S}_0'$.

The two candidate grids should not be chosen by the same rule, and the contrast between them is instructive. As \cref{\refsec:PEELOG} sets out, the residual \eqref{\refeq:stateResidual:iota} reaches $\iota_t$ \emph{only} through the interpolated continuation policies, so it carries no curvature of its own: its error is the kink structure of those interpolants, and it is removed by taking $\mathcal{S}_0'=\mathcal{S}_0$, or any refinement containing $\mathcal{S}_0$'s nodes, rather than by refining. The residual \eqref{\refeq:stateResidual:s} behaves in the opposite way. There $\bm{B}_{t+1}\left(s_t,h^{t+1}\right)$ is genuinely nonlinear in $s_t$, and that curvature swamps the same kinks, so alignment buys nothing -- aligned and unaligned grids of equal size perform indistinguishably -- while refinement pays. In our calibration the relative error of the located $s_t$ across the state grid is $2\cdot 10^{-2}$ at $\mathcal{S}'=\mathcal{S}$, $5\cdot 10^{-4}$ at four times the nodes, and $7\cdot 10^{-5}$ at four times the nodes spaced logarithmically. The last refinement earns its place because $s_t$ ranges over a factor of about twenty across the grid, so a uniform $\mathcal{S}'$ spends cells of constant width at the bottom of that range, where they are worth least. We therefore take $\mathcal{S}'$ finer than $\mathcal{S}$ and logarithmically spaced, and $\mathcal{S}_0'=\mathcal{S}_0$.

\smalltitle{The forward pass sees neither predetermined state.}
Suppose period $t+1$ is solved and we hold the continuation policies $\tau^{t+1}(\cdot,\cdot)$ and $h^{t+1}(\cdot,\cdot)$, both functions of the pair $(s_t,\iota_t)$ entering $t+1$. Here we carry $h_{t+1}$ itself rather than $\Theta_{h,t+1}$: with $s_t$ among the states, the level is a genuine function of the state, and $\bm{B}_{t+1}$ needs it as a level. For a candidate triple $(\tau_t, s_t, \iota_t)$ the auxiliary equations \eqref{\refeq:auxiliary} evaluate in one pass:
\begin{subequations}\label{\refeq:stateApprox}
\begin{align}
    \bm{B}_{t+1} &= \bm{B}_{t+1}\left(s_t,\, h^{t+1}(s_t,\iota_t)\right), & B_{t+1}^0 &= B_{t+1}^0\left(s_t,\, h^{t+1}(s_t,\iota_t)\right), \label{\refeq:stateApprox:B}\\
    \Gamma_{s,t} &= \Gamma_{s,t}\left(\bm{B}_{t+1},\, \tau^{t+1}(s_t,\iota_t)\right), &
    \Theta_{h,t} &= \Theta_{h,t}\left(\tau_t,\, \tau^{t+1}(s_t,\iota_t),\, \theta_{t+1},\, \Gamma_{s,t}\right), \label{\refeq:stateApprox:Gammas}\\
    \Theta_{s,t} &= \Theta_{s,t}\left(\Theta_{h,t},\, \Gamma_{s,t}\right). \label{\refeq:stateApprox:Thetas}
\end{align}
\end{subequations}
Two conditions close the period. The candidate savings level must reproduce itself through \eqref{\refeq:auxiliary:s}, and the candidate ratio through \eqref{\refeq:auxiliary:s0_s}:
\begin{subequations}\label{\refeq:stateResidual}
\begin{align}
    \mathcal{R}_t\left(\tau_t,s_t,\iota_t;\, s_{t-1}\right) &\equiv \Theta_{s,t}\left(\tau_t,s_t,\iota_t\right)\left(\frac{s_{t-1}}{\nu_t}\right)^{\frac{\alpha(1+\xi)}{1+\alpha\xi}}-s_t, \label{\refeq:stateResidual:s}\\
    \mathcal{R}_t^0\left(\tau_t,s_t,\iota_t\right) &\equiv \iota_t\left(B_{t+1}^0,\, \Theta_{s,t}\left(\tau_t,s_t,\iota_t\right),\, \tau^{t+1}(s_t,\iota_t)\right)-\iota_t. \label{\refeq:stateResidual:iota}
\end{align}
\end{subequations}
Nothing in \eqref{\refeq:stateApprox} depends on either predetermined state. $s_{t-1}$ enters \eqref{\refeq:stateResidual} only through the explicit factor in \eqref{\refeq:stateResidual:s}, and $\iota_{t-1}$ enters neither residual at all -- the latter for the same reason as in \cref{\refsec:PEELOG}, namely that $\Theta_{s,t}$ is a coefficient function from which the level of past savings has already been divided out. The expensive part of the state approximation is therefore a function of $(\tau_t,s_t,\iota_t)$ alone, evaluated once on $\mathcal{T}\times\mathcal{S}'\times\mathcal{S}_0'$: a three-dimensional object where a naive reading of the problem would suggest five.

\smalltitle{The two states unnest into two one-dimensional roots.}
The pair \eqref{\refeq:stateResidual} looks like a simultaneous root in $(s_t,\iota_t)$, which would be an uncomfortable two-dimensional problem to solve robustly at every $(\tau_t,s_{t-1})$. It is not, and the reason is that \eqref{\refeq:stateResidual:iota} involves neither predetermined state. We may therefore solve it \emph{first}, holding $(\tau_t,s_t)$ fixed:
\begin{align}\label{\refeq:iotaOfTauS}
    \iota_t = \iota_t\left(\tau_t,s_t\right) \quad\text{ as the root of }\quad \mathcal{R}_t^0\left(\tau_t,s_t,\cdot\right)\text{ along } \mathcal{S}_0',
\end{align}
substitute into $\Theta_{s,t}$ to obtain a function of $(\tau_t,s_t)$ alone, and then solve \eqref{\refeq:stateResidual:s} as a one-dimensional root in $s_t$ along $\mathcal{S}'$, broadcasting over $s_{t-1}\in\mathcal{S}$. The nesting is \emph{exact}, not an approximation: $\iota_t$ is fully determined by $(\tau_t,s_t)$ before $s_{t-1}$ is ever consulted, so no information is lost by eliminating it first. Two robust one-dimensional searches thus replace one fragile two-dimensional one, and the ordering is forced -- the reverse nesting is unavailable, since $s_t$ cannot be resolved without $s_{t-1}$.

\smalltitle{Feasibility.}
Three conditions must hold before a cell of $\mathcal{T}\times\mathcal{S}\times\mathcal{S}_0$ can enter the selection step.
\begin{enumerate}
    \item \emph{The ratio root exists.} \eqref{\refeq:iotaOfTauS} need not have a root inside $\mathcal{S}_0'$. Mark the $(\tau_t,s_t)$ pairs for which one was located.
    \item \emph{The savings root exists.} Large $\tau_t$ depresses savings enough that the implied $s_t$ falls below $l_s$, and conversely at small $\tau_t$; outside that window no equilibrium is located and every object downstream of $s_t$ is undefined. Mark the $(\tau_t,s_{t-1})$ pairs for which a root was found, and require at least two feasible $\tau_t$ per state before attempting to select a maximum. Since $s_t(\tau_t)$ is decreasing in $\tau_t$, the feasible set is expected to be a contiguous interval in $\tau_t$, but we do not rely on that: masking cells individually is no harder and does not silently mis-handle a non-monotone case.
    \item \emph{The consumption levels are positive.} This condition has no counterpart in the log case, where only log-derivatives were needed and levels never appeared. Here every term of \eqref{\refeq:PEEcrraTerms} carries a factor raised to the power $1-1/\rho$, which is undefined at a non-positive base. Of the four, $\tilde{c}_{1,t}^i$ and $\tilde{c}_{1,t}^0$ are sums of positive terms whenever $\tau_{t+1}\geq 0$ and $\theta_{t+1}\leq 1$, and $c_{2,t}^0$ is positive by $l_{\iota}>0$; only $c_{2,t}^i$ requires a check, since its bracket $s_{t-1,i}/s_{t-1}+A\left(1+\theta_t\left[\ldots\right]\right)$ can turn negative for a low-productivity type when $s_{t-1,i}/s_{t-1}$ is small. In implementation it is nevertheless cheapest to test all four rather than to rely on the argument, and to mark a cell infeasible when \emph{any} type's level is non-positive: $z_t$ sums over types, so one bad type is enough to spoil the sum, and testing catches the resulting non-integer power of a negative base at its source instead of letting it propagate.
\end{enumerate}
Conditions 2 and 3 live on $\mathcal{T}\times\mathcal{S}$ and condition 1 collapses onto it once $s_t$ is resolved, so the combined mask is two-dimensional: it never depends on $\iota_{t-1}$. That is worth stating because it is what keeps the selection step from having to be repeated per $\iota$-node with a different feasible sub-grid each time, and it is a consequence of $l_{\iota}>0$ having disposed of the only way $\iota_{t-1}$ could render a cell infeasible.

\smalltitle{Which state combinations are worth solving at.}
The two states are not independent. Along any equilibrium path $(s_t,\iota_t)$ is jointly produced by period $t$'s equilibrium at $(s_{t-1},\iota_{t-1})$ and the tax selected there, so the product grid $\mathcal{S}\times\mathcal{S}_0$ generally contains pairs that no path can visit. Spending evaluations on them is the lesser problem; the real hazard is that the interpolants $\tau^t(\cdot,\cdot)$ and $h^t(\cdot,\cdot)$ are built over the whole rectangle, so a call at an unreachable pair returns a perfectly finite number with no equilibrium content behind it, and nothing downstream marks it as such.

We therefore record, alongside the solution at each $t$, the \emph{reachable set}
\begin{align}\label{\refeq:reachable}
    \mathcal{P}_t \equiv \left\lbrace\left(s_t,\iota_t\right)\text{ implied by } \tau_t\left(s_{t-1},\iota_{t-1}\right) \mbox{ }:\mbox{ } \left(s_{t-1},\iota_{t-1}\right)\in\mathcal{S}\times\mathcal{S}_0\right\rbrace,
\end{align}
the set of states period $t$ can actually hand to $t+1$, and make two checks against it. First, $\mathcal{P}_t\subseteq\mathcal{S}\times\mathcal{S}_0$: if not, the grids are too narrow for the dynamics they are being asked to represent and should be widened, never clipped. Second, the simulated path of \cref{\refsec:PEEpath} must remain inside $\mathcal{P}_t$ at every $t$; a path that leaves it is being driven by extrapolated policy, and the result should be reported as unreliable rather than returned. Both are diagnostics on the solution, not constraints imposed during it: we keep the full product grid for simplicity, and let the containment check rather than the grid resolution carry the accuracy claim.

\smalltitle{Consumption levels and the young generations' discount factor.}
Once $s_t(\tau_t,s_{t-1})$ and $\iota_t(\tau_t,s_t)$ are known, all four marginal indirect utilities take the form (level)$^{1-1/\rho}\times$(log-derivative):
\begin{subequations}\label{\refeq:PEEcrraTerms}
\begin{align}
    \frac{d\upsilon_{1,t}^i}{d\tau_t} &= \left(\hat{c}_{1,t}^i\right)^{1-1/\rho}\frac{d\ln\left(\hat{c}_{1,t}^i\right)}{d\tau_t}, &
    \frac{d\upsilon_{1,t}^0}{d\tau_t} &= \left(\hat{c}_{1,t}^0\right)^{1-1/\rho}\frac{d\ln\left(\hat{c}_{1,t}^0\right)}{d\tau_t}, \\
    \frac{d\upsilon_{2,t}^i}{d\tau_t} &= \left(c_{2,t}^i\right)^{1-1/\rho}\frac{d\ln\left(c_{2,t}^i\right)}{d\tau_t}, &
    \frac{d\upsilon_{2,t}^0}{d\tau_t} &= \left(c_{2,t}^0\right)^{1-1/\rho}\frac{d\ln\left(c_{2,t}^0\right)}{d\tau_t}.
\end{align}
\end{subequations}
For both young generations the discount factor $B_{t+1}^j$ is itself a function of $\tau_t$ through $s_t$, so the factor $\left(1+B_{t+1}^j\right)$ in $\upsilon_{1,t}^j$ cannot be held fixed while differentiating. It is convenient to absorb it into the consumption level, defining for $j = 0,1,...,J$
\begin{align}\label{\refeq:hatc1}
    \hat{c}_{1,t}^j \equiv \left(1+B_{t+1}^j\right)^{\frac{1}{1-1/\rho}}\tilde{c}_{1,t}^j \qquad\Longrightarrow\qquad \upsilon_{1,t}^j = \frac{\left(\hat{c}_{1,t}^j\right)^{1-1/\rho}}{1-1/\rho},
\end{align}
so that a single numerical derivative of $\ln(\hat{c}_{1,t}^j)$ picks up both channels at once.\footnote{In implementation $\hat{c}_{1,t}^j$ is never formed as a level. Only two things are ever needed from it, and both avoid the $1/(1-1/\rho)$ power: $\left(\hat{c}_{1,t}^j\right)^{1-1/\rho} = \left(1+B_{t+1}^j\right)\left(\tilde{c}_{1,t}^j\right)^{1-1/\rho}$ and $\ln\left(\hat{c}_{1,t}^j\right) = \ln\left(1+B_{t+1}^j\right)/(1-1/\rho)+\ln\left(\tilde{c}_{1,t}^j\right)$. This matters as $\rho\rightarrow 1$: the literal definition raises $1+B_{t+1}^j$ to a power that diverges, overflowing in double precision well before $\rho$ reaches $1$ (at $\rho = 1.001$ the exponent is already $1001$), whereas neither expression above does.} The device is needed for the informal young as well as the formal, since they too save: $\upsilon_{1,t}^0$ has the same structure as $\upsilon_{1,t}^i$, with $B_{t+1}^0$ built on the informal return $R_{t+1}^0$ of \eqref{\refeq:auxiliary:R0}.

\smalltitle{Which derivatives may be taken on the grid.}
Three of the log-derivatives in \eqref{\refeq:PEEcrraTerms} --- those of $h_t$, $\hat{c}_{1,t}^i$ and $\hat{c}_{1,t}^0$ --- are approximated numerically along the $\tau_t$ dimension of the solution grid, separately for each $s_{t-1}$, since each depends on $\tau_t$ through the continuation policies and has no closed form. They are taken in the coordinate $x = \ln(1-\tau_t)$ of \cref{\refsec:PEELOG} rather than in $\tau_t$ itself, for the reason given there and with the same order-of-magnitude gain: the factor $1/(1-\tau_t)$ common to all three is affine in $x$ and is then reproduced exactly. The remaining two do \emph{not} follow that route:
\begin{itemize}
    \item $d\ln(c_{2,t}^i)/d\tau_t$ \textbf{must} be taken from its closed form, given in \eqref{\refmodeleq:PEE}. Differentiating it numerically along the grid would be wrong, not merely imprecise: $s_{t-1,i}/s_{t-1}$ varies along that grid, and the policy maker takes it as predetermined, so the grid derivative would include a channel that does not belong in the first order condition.
    \item $d\ln(c_{2,t}^0)/d\tau_t$ \emph{may} be taken either way, and for the reason set out in \cref{\refsec:PEELOG}: $\iota_{t-1}$ is a grid coordinate, held fixed by construction as $\tau_t$ varies, so the hazard above does not arise. We use the closed form \eqref{\refeq:dv20} for consistency and because it is cheaper.
\end{itemize}
Both closed forms are exactly the expressions already used in the log case, \eqref{\refmodeleq:PEELOG}, with one substitution: the closed-form $\partial\ln(\Theta_{h,t})/\partial\tau_t$ is replaced by the numerically obtained $d\ln(h_t)/d\tau_t$. Since $h_t = \Theta_{h,t}(s_{t-1}/\nu_t)^{\alpha\xi/(1+\alpha\xi)}$ and $s_{t-1}$ is held fixed under $\partial/\partial\tau_t$, the two coincide, and the log and CRRA cases share one formula rather than two.

\smalltitle{Numerical stability.}
Four points deserve attention beyond the feasibility conditions above.
\begin{itemize}
    \item \emph{$\rho = 1$ is not representable, except in the terminal period.} The powers in \eqref{\refeq:PEEcrraTerms} and the exponent in \eqref{\refeq:hatc1} both divide by $1-1/\rho$, so the recursion should refuse $\rho = 1$ outright and direct the caller to \cref{\refsec:PEELOG} rather than approximate it. The terminal period is a genuine exception and worth keeping available: it never forms $\hat{c}_{1,T}^j$, so it remains well defined at $\rho = 1$, where $\bm{B}_T$ collapses to the primitive $\bm{\beta}$ and every level factor to unity. It then coincides \emph{exactly}, not merely in the limit, with the terminal period of \cref{\refsec:PEELOG} --- state-independent in $s_{T-1}$ and equal to it to machine precision. That identity is the sharpest available test of the CRRA implementation, which is reason enough not to place the refusal where it would be out of reach.
    \item \emph{Scale disparity in $z_t$.} When $\rho$ is far from $1$, the level factors $(\cdot)^{1-1/\rho}$ differ by orders of magnitude across household types, so $z_t$ is a sum of very unequally scaled terms. The located crossing should not be trusted on its own: evaluate the residual at the selected $\tau_t$ and report it, so that a crossing produced by cancellation between two large terms is visible.
    \item \emph{Two interpolated surfaces per derivative.} Each of the three numerical derivatives composes $\tau^{t+1}(\cdot,\cdot)$ and $h^{t+1}(\cdot,\cdot)$, both of which are two-dimensional interpolants of previously selected policies. Kinks inherited from them are amplified by differentiation, so a light smoothing of $\tau_t(s_{t-1},\iota_{t-1})$ and $h_t(s_{t-1},\iota_{t-1})$ before interpolation matters more here than in the log case, where only one surface was involved. Two qualifications from \cref{\refsec:PEELOG} carry over unchanged and are worth repeating because they are easy to get backwards: the smoothed policy must be clipped back into $[l,u]$, and the smoothing belongs to the policy surfaces alone --- the interpolant through which the derivatives themselves are taken should be an interpolating one, since smoothing it degrades the derivative rather than improving it. With two states the smoothing is applied separably, along $\mathcal{S}$ and then along $\mathcal{S}_0$, in each case in the coordinate the grid is spaced in. The requirement that the smoother's knots be \emph{fixed} rather than selected from the data carries over unchanged and matters more here, since the calibration's outer problem differentiates through two interpolated surfaces per derivative rather than one; \cref{\refsec:PEELOG} gives the argument and the failure it prevents.

    \item \emph{The reporting step, not the search, is the memory bottleneck.} Steps 1--3 evaluate arrays of size $M\cdot M_s'\cdot M_{\iota}'$ and $M\cdot M_s$. Step 4 re-solves the state approximation at the selected taxes, of which there is one per state pair, so its largest intermediate is of size $M_s\cdot M_{\iota}\cdot M_s'\cdot M_{\iota}'$ --- the one object in the algorithm that scales with all four state and candidate grids at once, and hence as the fourth power of a uniform refinement. It should be evaluated in chunks over the state pairs rather than in a single pass. This is the one place where the remark closing \cref{\refalg:CRRA:grid} --- that cost is driven by the number of calls rather than the number of gridpoints --- does not apply.
    \item \emph{Extrapolation, not clamping.} A candidate $(s_t,\iota_t)$ from \eqref{\refeq:stateApprox} need not land on the nodes of $\mathcal{S}\times\mathcal{S}_0$, and clamping it to the boundary would silently return the boundary policy for an interior question. The interpolants should extrapolate, and the reachable-set check of \eqref{\refeq:reachable} is what bounds how far that extrapolation is trusted.
\end{itemize}

\begin{alg}[Backward grid search over the savings level and the informal savings ratio]
\label{\refalg:CRRA:grid}
Fix $\bm{\theta},\bm{\epsilon}$ and the grids $\mathcal{T}$, $\mathcal{S}$, $\mathcal{S}'$, $\mathcal{S}_0$, $\mathcal{S}_0'$. The recursion identifies a sequence of policy functions $\tau_t(s_{t-1},\iota_{t-1})$ together with the state policy functions $s_t(\tau_t,s_{t-1})$ and $\iota_t(\tau_t,s_t)$, iterating backwards from $t=T$.

\smallskip\noindent\emph{Terminal period.} There is no continuation policy, hence no fixed point: by \cref{\refmodelsec:finitehorizon} every object entering $z_T$ is closed form in $(\tau_T, s_{T-1}, \iota_{T-1})$. On the full grid $\mathcal{T}\times\mathcal{S}\times\mathcal{S}_0$ evaluate and store $\Theta_{h,T}(\tau_T)$ from \eqref{\refeq:auxiliary:ThetahT}, $d\ln(h_T)/d\tau_T = -\frac{\xi}{1+\alpha\xi}\frac{1}{1-\tau_T}$, $h_T$, $\bm{B}_T$, $\Gamma_{s,T-1}$, $s_{T-1,i}/s_{T-1}$, and the levels $\tilde{c}_{1,T}^i$, $c_{2,T}^i$ and $c_{2,T}^0$. The young terms are closed form by \eqref{\refmodeleq:terminalPEECRRA}, with $d\upsilon_{1,T}^0/d\tau_T = 0$ -- the informal young neither save nor face a future at $T$ -- so no numerical differentiation is needed in the terminal period. Assemble $z_T$, apply the selection rule \eqref{\refeq:candidates} along $\tau$ for each state pair, and obtain $\tau_T(s_{T-1},\iota_{T-1})$. Re-evaluate at the selected taxes and construct the interpolants $\tau^T(\cdot,\cdot)$ and $h^T(\cdot,\cdot)$.

The terminal period is the one place where the full three-dimensional grid is evaluated outright. That is affordable precisely because there is no state approximation to perform: every object is closed form, so the cost is one vectorized pass rather than a root-finding problem per cell.

\smallskip\noindent\emph{Recursion, $t = T-1,...,1$.} Given $\tau^{t+1}(\cdot,\cdot)$ and $h^{t+1}(\cdot,\cdot)$, each period proceeds in four steps.
\begin{enumerate}
    \item \emph{State approximation.} On $\mathcal{T}\times\mathcal{S}'\times\mathcal{S}_0'$ evaluate the forward pass \eqref{\refeq:stateApprox}. Locate the root of \eqref{\refeq:stateResidual:iota} along $\mathcal{S}_0'$ to obtain $\iota_t(\tau_t,s_t)$ on $\mathcal{T}\times\mathcal{S}'$, then form \eqref{\refeq:stateResidual:s} for every $s_{t-1}\in\mathcal{S}$ by broadcasting and locate its root along $\mathcal{S}'$, giving $s_t(\tau_t,s_{t-1})$ and, by composition, $\iota_t(\tau_t,s_{t-1})$ on $\mathcal{T}\times\mathcal{S}$ together with feasibility conditions 1 and 2. This is the only step that touches $\mathcal{S}'$ or $\mathcal{S}_0'$, and the only one that calls the continuation interpolants.
    \item \emph{Current-period objects.} With $s_t$ and $\iota_t$ resolved, evaluate on $\mathcal{T}\times\mathcal{S}$: $h_t\left(\Theta_{h,t},s_{t-1}\right)$, then $R_t$, $\bm{B}_t$, $\Gamma_{s,t-1}\left(\bm{B}_t,\tau_t\right)$ and $s_{t-1,i}/s_{t-1}$ at the previous period's vintage. Then the levels entering \eqref{\refeq:PEEcrraTerms}: $\hat{c}_{1,t}^i$ and $\hat{c}_{1,t}^0$ from \eqref{\refeq:hatc1}, and $c_{2,t}^i$. Apply feasibility condition 3.
    \item \emph{Derivatives and the first order condition.} Approximate $d\ln(h_t)/d\tau_t$, $d\ln(\hat{c}_{1,t}^i)/d\tau_t$ and $d\ln(\hat{c}_{1,t}^0)/d\tau_t$ numerically along $\tau_t$, separately for each $s_{t-1}$; take the two old-generation log-derivatives from their closed forms. Assemble the three state-independent terms of $z_t$ on $\mathcal{T}\times\mathcal{S}$, leaving infeasible cells undefined, then extend to $\mathcal{T}\times\mathcal{S}\times\mathcal{S}_0$ by evaluating $c_{2,t}^0$ and \eqref{\refeq:dv20} at each $\iota_{t-1}$ -- a broadcast, since neither requires re-entering \eqref{\refeq:stateApprox}.
    \item \emph{Selection and reporting.} Apply \eqref{\refeq:candidates} along $\tau_t$ for each state pair, skipping infeasible cells, to obtain $\tau_t(s_{t-1},\iota_{t-1})$. Re-evaluate steps 1--2 at the selected taxes to obtain the solution on $\mathcal{S}\times\mathcal{S}_0$, record the reachable set \eqref{\refeq:reachable}, and construct the smoothed interpolants $\tau^t(\cdot,\cdot)$ and $h^t(\cdot,\cdot)$ that period $t-1$ will call. The selected $s_t$, $\iota_t$ and $\Gamma_{s,t}$ may be recorded as interpolants over the state pair as well. Period $t-1$ never calls them: they serve only the cheaper variant of the path solve in \cref{\refsec:PEEpath}, which walks them forward instead of re-solving the transitions, and the accuracy it gives up in doing so is quantified there.
\end{enumerate}
Steps 1--4 are stated for the whole state grid at once rather than for one $(s_{t-1},\iota_{t-1})$ at a time. Nothing forces this --- the states are independent and could equally be looped over --- but the grid evaluations dominate the cost only through the number of \emph{calls}, not the number of gridpoints, so handling all states in one pass is materially cheaper than looping, at the price of carrying the feasibility mask explicitly rather than subsetting the grid per state. Step 4 is the exception noted above: its re-solve carries an intermediate that scales with all four state and candidate grids simultaneously, and should be chunked over the state pairs.
\end{alg}
